\section{The \instname\ data reduction flow.}
\label{sec:reduction_flow}


The overall data flow of the \instname\ pipeline is displayed in Figure \ref{fig:reduction_cascade}.

 
\begin{figure}[h!]
\centering
\includegraphics[width=0.9\textwidth]{../instrument_tutorial/figures/reduction_cascade.jpg}
\caption{The data reduction cascade of the \instname\ workflow.}
\label{fig:reduction_cascade}
\end{figure}



The reduction cascade is organized in tasks, which represent
well-defined steps in the process. Tasks can be grouped inside
sub-workflows.  Each task runs a recipe; the detailed description of
the algorithms, input, outputs and recipe parameters used in each
recipe are available in the pipeline manual. Here, we present only the
description of most important features.

The \wkfname\  \edps\ workflow is designed to execute the tasks that
deliver the final reduced data cube for each dataset. It can be either
the product of a single exposure, or the combination of multiple
exposures. Only calibrations needed by the selected the scientific
exposures are processed.

It is possible to set \edps\ to perform the data reduction until a
certain step of the reduction chain (e.g. to reduce only standar
stars, or only flat fields).  This is done by specifying the desired
tasks in the field {\bf Select reduction target} of the {\bf Raw Data}
tab.

The reduction steps of the \wkfname\ workflow are listed below. Before
starting the reduction, the parameters of the recipes associated to
each task can be configured by pressing the button

\includegraphics[width=0.8cm]{../common/figures/configure_dataset.jpg}
close to each dataset configuration.
See for more info on the configuration editor \ref{sec:configuration_editor}





\subsection{Process dark exposure}

This step is carried out by the task **dark** using the pipeline recipe {\tt kmos\_dark}. 
It produces a MASTER\_DARK and a BADPIXEL\_DARK frame. Per default, only the bad pixel map from this step
is used in the reduction chain.

\subsection{Create Master Flat}

Task: {\bf lamp\_flat}. Recipe: {\tt kmos\_flat}. 

This step uses exposures with the flat-field lamp to create a MASTER\_FLAT, files that describe the geometry 
of the IFU slitlets on the detector (XCAL, YCAL, and SLIT\_EDGE), and a BADPIXEL\_FLAT frame.

\subsection{Wavelength Calibration}

Task: {\bf arc}. Recipe: {\tt kmos\_wave\_cal}. 

This step uses exposures with the Argon and Neon arc lamps to create the wavelength calibration LCAL.
A useful product for verifying the wavelength calibration is DET\_IMG\_WAVE which contains the wavelength-calibrated
images of the input frames.

\subsection{Illumination Correction}

Task: {\bf illumination}. Recipe: {\tt kmos\_illumination}. 

This step can be run on lamp-flat or sky-flat exposures. It creates
for each IFU an image (product: ILLUM\_CORR) for correcting the
sensitivity variation within the IFU. It is recommended to execute
this step on lamp flats (which is the default workflow setup).

This step can be customized by setting the workflow parameter
`use\_sky\_flats` to 'false' (default, recommended) or 'true'.

\subsection{Standard star}

Task: {\bf standard}. Recipe: {\tt kmos\_std\_star}. 

This step processes a telluric standard exposure and extracts a spectrum for each IFU that was used.
Products:
\begin{itemize}
\item STAR\_SPEC: extracted spectrum
\item  TELLURIC: normalized telluric spectrum
\item  STD\_IMAGE: reconstructed IFU field-of-view image
\item  STD\_MASK: mask for extracting the standard star.
\end{itemize}

The further usage of the products is controlled via the workflow parameter `molecfit` which can have the values:

\begin{itemize}
\item  'standard': Atmospheric features are modeled with the molecfit algorithm using the telluric standard as
                reference. The results are used to construct the full atmospheric transmission to correct
                science exposures. A static response curve is used to correct the relative instrumental 
                efficiency with wavelength. The zeropoint computed on the observed telluric standard is used for
                absolute spectrophotometric calibration.
\item  'science': Atmospheric features are modeled with the molecfit algorithm using the scientific exposure itself
                as reference. To be used only if the science exposure has a bright continuum that allows
                a proper fit to the atmospheric features. A static response curve is used to correct the
                relative instrumental efficiency with wavelength. The zeropoint computed
                on the observed telluric standard is used for absolute spectrophotometric calibration.
\item  'false': The observations of the telluric standard stars are compared with a model spectrum for that star.
                This generates a combined correction that includes response curve and telluric features, which is then
                used to correct scientific exposures.
\end{itemize}

\subsection{Response}

This step involves either the task {\bf response} or the task {\bf telluric\_and\_response}. Both use the pipeline recipe
{\tt kmos\_gen\_telluric}.

If the workflow parameter `molecfit` is set to 'false', then the task **telluric\_and\_response** is used. 
The telluric spectrum as calculated from the {\bf standard} task and the response are passed to the science reduction.
For this other values of the workflow parameter, only the response is used.

\subsection{Science reduction}

Task: {\bf object}. Recipe: {\tt kmos\_sci\_red}.

In this step, the basic data reduction of the science data is executed. The data from the IFUs are extracted, 
calibrated, and reconstructed into data cubes. If the workflow parameter `molecfit` is set to 'false', then the 
products are flux-calibrated and corrected for telluric absorption. Otherwise, the products are only
flux-calibrated. The telluric correction is then executed in the following workflow steps.

Main products:
\begin{itemize}

\item SCI\_RECONSTRUCTED: one file for each exposure that contains
  IFUs on objects (no file for pure sky exposures). Each file has 24
  extensions with data cubes.
  
\item SCI\_RECONSTRUCTED\_COLL: files with 24 extensions containing
  the reconstructed field-of-view images for each IFU.

\end{itemize}

\subsection{Telluric model}

This step either involves the tasks {\bf model\_on\_standard} and {\bf transmission\_on\_standard}
or the tasks {\bf model\_on\_science} and {\bf transmission\_on\_science}. The tasks are encapsulated in the
sub workflows telluric\_on\_standard and telluric\_on science, respectively.
The pipeline recipes {\tt kmos\_molecfit\_model},  {\tt kmos\_molecfit\_calctrans} are used.

If the workflow parameter `molecfit` is set to 'standard', then 
the atmospheric model is fit to the extracted spectra from the {\bf reponse} task. 
With the parameter set to `science`, the fit is executed on the science data themselves.
In both cases, the actual transmission curve is calculated using the model and the information from the science data.

\subsection{Telluric correction}

Task: {\bf telluric\_correction}. Recipe: {\tt kmos\_molecfit\_correct}.

In this step, the science data are corrected for telluric absorption using the transmission from the previous step.

\subsection{Object combination}
\label{sec:object_combination}
Task: {\bf object\_combination}. Recipe: {\tt kmos\_combine}.

The different exposures for each object are combined. The final products are:
\begin{itemize}
\item One combined data cube for each object in the dataset, coming in two different formats as
  `COMBINED\_CUBE` and `IDP\_COMBINED\_CUBE`. Both products have the same data content. 
  The latter file has the needed metadata information for ingestion into the ESO archive.
\item A field-of-view image `COMBINED\_IMAGE` and an exposure mask image `EXP\_MASK` for each object.
\end{itemize}

\subsection{Cubes combination}
\label{sec:cubes_combination}
Task: {\bf cubes\_combination}. Recipe: {\tt kmos\_combine}.

If data have already been processed with the previous reduction then the resulting cubes can be combined
for each object. This is useful for combining observations of the same targets which have been measured
using different Observation Blocks.


